% topic_01.tex — Content only. Compiled via main.tex.
% ── No \documentclass, \begin{document} or \end{document} here ──

\chapteropener{1}{Physical Quantities and Units}{%
  \item Physical Quantities and Estimation
  \item SI Base Units and Derived Units
  \item Homogeneity of Equations
  \item SI Prefixes
  \item Errors, Uncertainties, Precision and Accuracy
  \item Scalars and Vectors
}

% ===== SECTION 1: PHYSICAL QUANTITIES =====
\section{Physical Quantities}

\begin{definitionbox}{Physical Quantity}
Every \textbf{physical quantity} consists of two parts:
$$\text{physical quantity} = \text{numerical magnitude} \times \text{unit}$$
For example: a length of $3.5\ \text{m}$ has magnitude $3.5$ and unit metres. A quantity without a unit (or with units that cancel) is \textbf{dimensionless}.
\end{definitionbox}

\begin{keybox}{Reasonable Estimates}
The syllabus requires you to make sensible estimates of physical quantities. Useful benchmarks to memorise:

\begin{center}
\renewcommand{\arraystretch}{1.4}
\begin{tabular}{@{} l l l @{}}
\toprule
\textbf{Quantity} & \textbf{Estimate} & \textbf{Notes} \\
\midrule
Mass of a person & $70\ \text{kg}$ & \\
Height of a person & $1.7\ \text{m}$ & \\
Mass of a car & $1000\ \text{kg}$ & \\
Speed of a car on motorway & $30\ \text{m s}^{-1}$ & ($\approx 110\ \text{km h}^{-1}$) \\
Speed of sound in air & $340\ \text{m s}^{-1}$ & \\
Speed of light & $3\times10^8\ \text{m s}^{-1}$ & \\
Atmospheric pressure & $1\times10^5\ \text{Pa}$ & \\
Density of water & $1000\ \text{kg m}^{-3}$ & \\
Density of air & $1.2\ \text{kg m}^{-3}$ & \\
Diameter of an atom & $\sim10^{-10}\ \text{m}$ & \\
Diameter of a nucleus & $\sim10^{-15}\ \text{m}$ & \\
\bottomrule
\end{tabular}
\end{center}
\end{keybox}


\pagebreak

% ===== SECTION 2: SI UNITS =====
\section{SI Base Units}

\begin{definitionbox}{SI Base Quantities and Units}
The International System of Units (SI) defines \textbf{seven base quantities}. The five required for this syllabus are:

\begin{center}
\renewcommand{\arraystretch}{1.5}
\begin{tabular}{@{} l l l @{}}
\toprule
\textbf{Base quantity} & \textbf{SI unit} & \textbf{Symbol} \\
\midrule
Mass & kilogram & kg \\
Length & metre & m \\
Time & second & s \\
Electric current & ampere & A \\
Temperature & kelvin & K \\
\bottomrule
\end{tabular}
\end{center}
\end{definitionbox}

\begin{definitionbox}{Derived Units}
All other units are \textbf{derived} from the base units by multiplication or division. Examples:

\begin{center}
\renewcommand{\arraystretch}{1.5}
\begin{tabular}{@{} l l l @{}}
\toprule
\textbf{Quantity} & \textbf{Derived unit} & \textbf{In base units} \\
\midrule
Force & newton (N) & kg m s$^{-2}$ \\
Energy / Work & joule (J) & kg m$^2$ s$^{-2}$ \\
Power & watt (W) & kg m$^2$ s$^{-3}$ \\
Pressure & pascal (Pa) & kg m$^{-1}$ s$^{-2}$ \\
Charge & coulomb (C) & A s \\
Potential difference & volt (V) & kg m$^2$ s$^{-3}$ A$^{-1}$ \\
Resistance & ohm ($\Omega$) & kg m$^2$ s$^{-3}$ A$^{-2}$ \\
Frequency & hertz (Hz) & s$^{-1}$ \\
\bottomrule
\end{tabular}
\end{center}
\end{definitionbox}

\pagebreak


% ===== SECTION 3: HOMOGENEITY =====
\section{Homogeneity of Physical Equations}

\begin{definitionbox}{Homogeneity}
A physical equation is \textbf{homogeneous} if the units (dimensions) on both sides are identical. Every valid physical equation must be homogeneous.

\textbf{Method}: express every quantity in SI base units, then check both sides match.
\end{definitionbox}

\begin{keybox}{Using Homogeneity to Check Equations}
\begin{enumerate}[itemsep=3pt]
  \item Replace each symbol with its SI base units.
  \item Simplify both sides independently.
  \item If both sides have the same base units $\Rightarrow$ equation is \textbf{homogeneous} (possibly correct).
  \item If the sides differ $\Rightarrow$ equation is \textbf{definitely wrong}.
\end{enumerate}
\textit{Important caveat}: homogeneity is a necessary but not sufficient condition. A homogeneous equation can still be wrong (e.g.\ a missing numerical factor of 2 would not be detected).
\end{keybox}

\begin{examplebox}{Worked Example — Homogeneity Check}
Check whether $v^2 = u^2 + 2as$ is homogeneous.

LHS: $[v^2] = (\text{m s}^{-1})^2 = \text{m}^2\ \text{s}^{-2}$

RHS: $[u^2] = \text{m}^2\ \text{s}^{-2}$;\quad $[2as] = (\text{m s}^{-2})(\text{m}) = \text{m}^2\ \text{s}^{-2}$

Both terms on the RHS and the LHS have units m$^2$ s$^{-2}$ $\Rightarrow$ equation is \textbf{homogeneous}. \checkmark
\end{examplebox}

\pagebreak

% ===== SECTION 4: PREFIXES =====
\section{SI Prefixes}

\begin{center}
\renewcommand{\arraystretch}{1.5}
\begin{tabular}{@{} l l l r @{}}
\toprule
\textbf{Prefix} & \textbf{Symbol} & \textbf{Multiplier} & \textbf{Power of 10} \\
\midrule
tera  & T & 1\,000\,000\,000\,000 & $10^{12}$ \\
giga  & G & 1\,000\,000\,000      & $10^{9}$  \\
mega  & M & 1\,000\,000           & $10^{6}$  \\
kilo  & k & 1\,000                & $10^{3}$  \\
deci  & d & 0.1                   & $10^{-1}$ \\
centi & c & 0.01                  & $10^{-2}$ \\
milli & m & 0.001                 & $10^{-3}$ \\
micro & $\mu$ & 0.000\,001        & $10^{-6}$ \\
nano  & n & 0.000\,000\,001       & $10^{-9}$ \\
pico  & p & 0.000\,000\,000\,001  & $10^{-12}$ \\
\bottomrule
\end{tabular}
\end{center}

\begin{warningbox}{Common Prefix Errors}
\begin{itemize}[itemsep=2pt]
  \item $1\ \text{nm} = 1\times10^{-9}\ \text{m}$, \textbf{not} $10^{-6}\ \text{m}$ — don't confuse nano (n) with micro ($\mu$).
  \item When squaring or cubing a unit with a prefix, the prefix is also raised to that power: $1\ \text{cm}^2 = (10^{-2}\ \text{m})^2 = 10^{-4}\ \text{m}^2$, not $10^{-2}\ \text{m}^2$.
  \item $1\ \text{mm}^3 = (10^{-3}\ \text{m})^3 = 10^{-9}\ \text{m}^3$.
\end{itemize}
\end{warningbox}

\pagebreak

% ===== SECTION 5: ERRORS AND UNCERTAINTIES =====
\section{Errors, Uncertainties, Precision and Accuracy}

\begin{definitionbox}{Types of Error}
\begin{itemize}[itemsep=4pt]
  \item \textbf{Systematic error}: an error that affects all readings by the same amount in the same direction. The mean of repeated readings is shifted from the true value. Cannot be reduced by repeating measurements.
  \begin{itemize}
    \item Examples: zero error on an instrument; wrongly calibrated scale; parallax error if always reading from the same angle.
  \end{itemize}
  \item \textbf{Random error}: an error that causes readings to scatter unpredictably above and below the true value. Can be reduced by taking more readings and averaging.
  \begin{itemize}
    \item Examples: reaction time variation; reading a scale to the nearest division; electrical noise.
  \end{itemize}
  \item \textbf{Zero error}: a specific systematic error where an instrument reads non-zero when it should read zero. Corrected by subtracting the zero reading from all measurements.
\end{itemize}
\end{definitionbox}

\begin{definitionbox}{Precision and Accuracy}
\begin{itemize}[itemsep=4pt]
  \item \textbf{Precision}: how close repeated measurements are to \emph{each other}. A precise instrument gives small random errors (small spread). Precision is about \textbf{repeatability}.
  \item \textbf{Accuracy}: how close a measurement is to the \emph{true value}. An accurate instrument has small systematic error. Accuracy is about \textbf{closeness to truth}.
\end{itemize}
It is possible to be precise but inaccurate (systematic error shifts all readings the same way) or accurate but imprecise (readings scattered around the true value).
\end{definitionbox}

\begin{center}
\begin{tikzpicture}[scale=1.5]
  % Four target diagrams
  \foreach \cx/\cy/\lab/\subl in {
    0/0/{Precise, accurate}/{small random, small systematic},
    4/0/{Precise, not accurate}/{small random, large systematic},
    0/-3.5/{Not precise, accurate}/{large random, small systematic},
    4/-3.5/{Not precise, not accurate}/{large random, large systematic}
  }{
    % Target rings
    \foreach \r in {0.9,0.6,0.3}{
      \draw[midblue, thin] (\cx,\cy) circle (\r);
    }
    % True value centre
    \fill[darkblue] (\cx,\cy) circle (2.5pt);
    \node[darkblue, font=\small, anchor=south] at (\cx, \cy+1.0) {\lab};
   \node[darkgray, font=\small, anchor=north] at (\cx, \cy-1.05) {\subl};
  }
  % Dots for precise & accurate
  \foreach \dx/\dy in {0.05/0.08,-0.07/0.04,0.04/-0.06,-0.03/0.09,0.06/-0.03}{
    \fill[pinkframe] (0+\dx, 0+\dy) circle (2pt);
  }
  % Dots for precise, not accurate
  \foreach \dx/\dy in {0.55/0.08,0.45/0.05,0.58/-0.04,0.52/0.09,0.47/-0.06}{
    \fill[pinkframe] (4+\dx, 0+\dy) circle (2pt);
  }
  % Dots for not precise, accurate (scattered around centre)
  \foreach \dx/\dy in {0.3/0.2,-0.4/0.1,0.1/-0.5,0.5/-0.3,-0.2/0.4,0.04/-.03}{
    \fill[pinkframe] (0+\dx, -3.5+\dy) circle (2pt);
  }
  % Dots for not precise, not accurate
  \foreach \dx/\dy in {0.5/0.4,-0.3/0.7,0.8/-0.3,-0.5/0.1,0.2/0.6}{
    \fill[pinkframe] (4+\dx, -3.5+\dy) circle (2pt);
  }
\end{tikzpicture}
\end{center}

\begin{formulabox}{Combining Uncertainties}
\begin{center}
\renewcommand{\arraystretch}{1.6}
\begin{tabular}{@{} l l @{}}
\toprule
\textbf{Operation} & \textbf{Rule for uncertainty} \\
\midrule
$z = x + y$ or $z = x - y$ & $\Delta z = \Delta x + \Delta y$ \quad (add \textbf{absolute} uncertainties) \\
$z = xy$ or $z = x/y$ & $\dfrac{\Delta z}{z} = \dfrac{\Delta x}{x} + \dfrac{\Delta y}{y}$ \quad (add \textbf{percentage} uncertainties) \\
$z = x^n$ & $\dfrac{\Delta z}{z} = \left|n\right| \dfrac{\Delta x}{x}$ \quad (multiply percentage uncertainty by $|n|$) \\
\bottomrule
\end{tabular}
\end{center}
\vspace{0.4em}
\textbf{Absolute uncertainty} $\Delta x$: has the same unit as $x$.\\
\textbf{Percentage uncertainty}: $\dfrac{\Delta x}{x} \times 100\%$ — dimensionless.
\end{formulabox}

\begin{warningbox}{Common Uncertainty Mistakes}
\begin{itemize}[itemsep=2pt]
  \item For \textbf{addition/subtraction}, always add \emph{absolute} uncertainties — never subtract.
  \item For \textbf{powers}: $z = x^2$ gives $\Delta z / z = 2\,\Delta x/x$ — the power doubles the percentage uncertainty.
  \item A constant with no uncertainty (e.g.\ $\pi$, $2$, $g$ if treated as exact) does not contribute to the uncertainty.
  \item When reading an analogue scale, the uncertainty is typically $\pm$ half the smallest division.
\end{itemize}
\end{warningbox}

\begin{examplebox}{Worked Example — Combining Uncertainties}
A cylinder has radius $r = (1.50 \pm 0.02)\ \text{cm}$ and height $h = (4.00 \pm 0.05)\ \text{cm}$. Calculate the volume and its percentage uncertainty.

$V = \pi r^2 h = \pi (1.50)^2(4.00) = 28.3\ \text{cm}^3$

Percentage uncertainties:
$$\frac{\Delta r}{r} = \frac{0.02}{1.50} = 1.33\%$$
$$\frac{\Delta h}{h} = \frac{0.05}{4.00} = 1.25\%$$
$$\frac{\Delta V}{V} = 2\times1.33\% + 1.25\% = 3.91\% \approx \mathbf{3.9\%}$$
($r$ is squared so its percentage uncertainty is doubled; $\pi$ contributes nothing.)

Absolute uncertainty: $\Delta V = 0.039 \times 28.3 = 1.1\ \text{cm}^3$

$\therefore V = (28.3 \pm 1.1)\ \text{cm}^3$
\end{examplebox}

% ===== SECTION 6: SCALARS AND VECTORS =====
\section{Scalars and Vectors}

\begin{definitionbox}{Scalars and Vectors}
\begin{itemize}[itemsep=4pt]
  \item A \textbf{scalar} quantity has \textbf{magnitude only}.
  \item A \textbf{vector} quantity has both \textbf{magnitude and direction}.
\end{itemize}
\begin{center}
\renewcommand{\arraystretch}{1.4}
\begin{tabular}{@{} l l @{}}
\toprule
\textbf{Scalars} & \textbf{Vectors} \\
\midrule
distance, speed & displacement, velocity \\
mass, density & acceleration, force, weight \\
energy, power, work & momentum, impulse \\
temperature, pressure & electric field strength \\
time, frequency & gravitational field strength \\
\bottomrule
\end{tabular}
\end{center}
\end{definitionbox}

\begin{formulabox}{Adding Vectors}
Vectors must be added \textbf{tip-to-tail}. The resultant $\vec{R} = \vec{A} + \vec{B}$ is the single vector drawn from the tail of $\vec{A}$ to the tip of $\vec{B}$.

For two vectors at an angle $\theta$ to each other, the magnitude of the resultant:
$$R = \sqrt{A^2 + B^2 + 2AB\cos\theta}$$
For \textbf{perpendicular} vectors ($\theta = 90^\circ$): $R = \sqrt{A^2 + B^2}$, direction $= \arctan(B/A)$.
\end{formulabox}

\begin{formulabox}{Subtracting Vectors}
To find $\vec{A} - \vec{B}$, reverse the direction of $\vec{B}$ to get $-\vec{B}$, then add tip-to-tail:
$$\vec{A} - \vec{B} = \vec{A} + (-\vec{B})$$
\begin{itemize}[itemsep=2pt]
  \item $-\vec{B}$ has the \textbf{same magnitude} as $\vec{B}$ but \textbf{opposite direction}.
  \item This is used whenever you need a \textbf{change in a vector quantity}, e.g.\ change in velocity $\Delta\vec{v} = \vec{v}_2 - \vec{v}_1$.
  \item For perpendicular vectors: $|\vec{A} - \vec{B}| = \sqrt{A^2 + B^2}$ (same magnitude as addition, different direction).
\end{itemize}
\end{formulabox}

\begin{center}
\begin{tikzpicture}[scale=1.05, every node/.style={font=\small}]
  % --- Vector addition ---
  \begin{scope}[xshift=0cm]
    \node[darkblue, font=\small] at (1.75, 2.6) {$\vec{A} + \vec{B}$};
    \draw[-latex, midblue, very thick] (0,0) -- (2.5,0);
    \node[midblue, anchor=north] at (1.25,-0.08) {$\vec{A}$};
    \draw[-latex, pinkframe, very thick] (2.5,0) -- (3.5,1.8);
    \node[pinkframe, anchor=west] at (3,0.9) {$\vec{B}$};
    \draw[-latex, darkblue, very thick] (0,0) -- (3.5,1.8);
    \node[darkblue, anchor=south east] at (1.55,1.05) {$\vec{R}$};
  \end{scope}

  % --- Vector subtraction ---
  \begin{scope}[xshift=5.5cm]
    \node[darkblue, font=\small] at (1.75, 2.6) {$\vec{A} - \vec{B}$};
    % A
    \draw[-latex, midblue, very thick] (0,0) -- (2.5,0);
    \node[midblue, anchor=north] at (1.25,-0.08) {$\vec{A}$};
    % Original B (dashed, for reference)
    \draw[-latex, pinkframe!35, dashed, thick] (2.5,0) -- (3.5,1.8);
    \node[pinkframe!60, anchor=west, font=\small] at (3,0.9) {$\vec{B}$};
    % -B: reversed from tip of A
    \draw[-latex, pinkframe, very thick] (2.5,0) -- (1.5,-1.8);
    \node[pinkframe, anchor=west] at (1.95,-0.95) {$-\vec{B}$};
    % Resultant of A + (-B)
    \draw[-latex, darkblue, very thick] (0,0) -- (1.5,-1.8);
    \node[darkblue, anchor=north east] at (0.65,-0.95) {$\vec{A}{-}\vec{B}$};
  \end{scope}

  % --- Vector triangle (equilibrium) ---
  \begin{scope}[xshift=11.0cm]
    \node[darkblue, font=\small] at (1.1, 2.6) {Equilibrium};
    \draw[-latex, midblue, very thick] (0,0) -- (2.2,0);
    \node[midblue, anchor=north] at (1.1,-0.08) {$\vec{F_1}$};
    \draw[-latex, pinkframe, very thick] (2.2,0) -- (2.2,1.7);
    \node[pinkframe, anchor=west] at (2.25,0.85) {$\vec{F_2}$};
    \draw[-latex, lightgray, very thick] (2.2,1.7) -- (0,0);
    \node[lightgray, anchor=south east] at (1.1,0.9) {$\vec{F_3}$};
    \node[darkblue, font=\small, anchor=north] at (1.1,-0.5) {closed triangle};
  \end{scope}
\end{tikzpicture}
\end{center}

\begin{keybox}{Change in a Vector Quantity}
A common application of vector subtraction is finding the \textbf{change in velocity} $\Delta\vec{v}$:
$$\Delta\vec{v} = \vec{v}_2 - \vec{v}_1 = \vec{v}_2 + (-\vec{v}_1)$$
Draw $\vec{v}_1$ and $\vec{v}_2$ from the same point. $\Delta\vec{v}$ is the vector from the tip of $\vec{v}_1$ to the tip of $\vec{v}_2$.

\textit{Example}: a ball bouncing off a wall at the same speed. Although $|\vec{v}_1| = |\vec{v}_2|$, the direction has changed so $\Delta\vec{v} \neq 0$ — there is a non-zero acceleration (and therefore a resultant force).
\end{keybox}

\begin{formulabox}{Resolving a Vector into Components}
Any vector $\vec{F}$ at angle $\theta$ to the horizontal can be split into two perpendicular components:
\begin{center}
\Large $F_x = F\cos\theta \qquad F_y = F\sin\theta$
\end{center}

\begin{center}
\begin{tikzpicture}[scale=1.2, every node/.style={font=\small}]
  % Vector
  \draw[-latex, darkblue, very thick] (0,0) -- (3.0,2.0);
 \node[darkgray, anchor=south west] at (1.4,1.1) {$F$};
  % Horizontal component
  \draw[-latex, midblue, thick] (0,0) -- (3.0,0);
  \node[midblue, anchor=north] at (1.5,-0.05) {$F\cos\theta$};
  % Vertical component
  \draw[-latex, pinkframe, thick] (3.0,0) -- (3.0,2.0);
  \node[pinkframe, anchor=west] at (3.05,1.0) {$F\sin\theta$};
  % Dashed lines
  \draw[dashed, darkgray] (0,2.0) -- (3.0,2.0);
  \draw[dashed, darkgray] (0,0) -- (0,2.0);
  % Angle arc
  \draw[darkblue] (0.7,0) arc (0:33.7:0.7);
  \node[darkblue] at (0.95,0.2) {$\theta$};
\end{tikzpicture}
\end{center}
\end{formulabox}

\begin{keybox}{Coplanar Forces in Equilibrium — Vector Triangle}
Three coplanar forces in equilibrium can be represented by a \textbf{closed triangle}: draw the vectors tip-to-tail; if they form a closed triangle (the last tip meets the first tail), the system is in equilibrium. This provides a graphical method for finding unknown force magnitudes or directions.
\end{keybox}

% ===== SECTION 7: FORMULA SUMMARY =====
\section{Formula Summary Sheet}

\begin{minipage}{\linewidth}
\begin{tcolorbox}[colback=lightgold, colframe=accentgold]
\renewcommand{\arraystretch}{1.8}
\begin{tabular}{@{}lll@{}}
\toprule
\textbf{Formula / Rule} & \textbf{Use} & \\
\midrule
$z = A + B \Rightarrow \Delta z = \Delta A + \Delta B$ & Add/subtract: add absolute uncertainties & \\
$z = AB \Rightarrow \Delta z/z = \Delta A/A + \Delta B/B$ & Multiply/divide: add percentage uncertainties & \\
$z = A^n \Rightarrow \Delta z/z = |n|\,\Delta A/A$ & Powers: multiply \% uncertainty by $|n|$ & \\
$F_x = F\cos\theta$ & Horizontal component of vector & \\
$F_y = F\sin\theta$ & Vertical component of vector & \\
$R = \sqrt{A^2+B^2}$ & Resultant of two perpendicular vectors & \\
\bottomrule
\end{tabular}
\end{tcolorbox}

\vspace{0.5em}
\begin{tcolorbox}[colback=lightblue, colframe=darkblue]
\textbf{Key facts:}\\[0.3em]
SI base units: kg, m, s, A, K.\\[0.2em]
\textbf{Systematic error}: shifts all readings same way; not reduced by repeating.\\[0.2em]
\textbf{Random error}: causes scatter; reduced by averaging more readings.\\[0.2em]
\textbf{Precision}: repeatability (small spread). \textbf{Accuracy}: closeness to true value.\\[0.2em]
\textbf{Homogeneity}: units must match on both sides of any valid equation.
\end{tcolorbox}
\end{minipage}

\pagebreak

% ===== SECTION 8: WORKED EXAMPLES =====
\section{Worked Examples}

\subsection*{Example 1 --- Derived Units in Base Units}

\textbf{Question:} Show that the unit of pressure (Pa) is equivalent to kg m$^{-1}$ s$^{-2}$.

\begin{examplebox}{Solution}
Pressure $= $ Force / Area. Force has units N $=$ kg m s$^{-2}$. Area has units m$^2$.
$$[\text{Pa}] = \frac{\text{N}}{\text{m}^2} = \frac{\text{kg m s}^{-2}}{\text{m}^2} = \text{kg m}^{-1}\ \text{s}^{-2} \quad \checkmark$$
\end{examplebox}

\subsection*{Example 2 --- Homogeneity Check}

\textbf{Question:} The period of a simple pendulum is given by $T = 2\pi\sqrt{L/g}$. Check this equation for homogeneity.

\begin{examplebox}{Solution}
LHS: $[T] = \text{s}$

RHS: $2\pi$ is dimensionless. $[L/g] = \text{m} / (\text{m s}^{-2}) = \text{s}^2$

$[\sqrt{L/g}] = \sqrt{\text{s}^2} = \text{s}$

Both sides have units of seconds $\Rightarrow$ equation is \textbf{homogeneous}. \checkmark
\end{examplebox}

\subsection*{Example 3 --- Combining Uncertainties}

\textbf{Question:} The resistance of a wire is found using $R = V/I$, where $V = (6.0 \pm 0.2)\ \text{V}$ and $I = (0.30 \pm 0.01)\ \text{A}$. Calculate $R$ and its percentage uncertainty.

\begin{examplebox}{Solution}
$R = V/I = 6.0/0.30 = \mathbf{20\ \Omega}$

Since this is division, add percentage uncertainties:
$$\frac{\Delta R}{R} = \frac{\Delta V}{V} + \frac{\Delta I}{I} = \frac{0.2}{6.0} + \frac{0.01}{0.30} = 3.33\% + 3.33\% = \mathbf{6.7\%}$$

Absolute uncertainty: $\Delta R = 0.067 \times 20 = 1.3\ \Omega$

$R = (20 \pm 1)\ \Omega$
\end{examplebox}

\subsection*{Example 4 --- Resolving and Adding Vectors}

\textbf{Question:} Two forces act on a point: $F_1 = 8.0\ \text{N}$ horizontally and $F_2 = 6.0\ \text{N}$ vertically upward. Find the magnitude and direction of the resultant.

\begin{examplebox}{Solution}
Since the forces are perpendicular:
$$R = \sqrt{F_1^2 + F_2^2} = \sqrt{8.0^2 + 6.0^2} = \sqrt{64 + 36} = \sqrt{100} = \mathbf{10\ \text{N}}$$

Direction above the horizontal:
$$\theta = \arctan\!\left(\frac{F_2}{F_1}\right) = \arctan\!\left(\frac{6.0}{8.0}\right) = \arctan(0.75) = \mathbf{36.9^\circ}$$

Resultant: $10\ \text{N}$ at $36.9^\circ$ above the horizontal.
\end{examplebox}

\subsection*{Example 5 --- Vector Subtraction (Change in Velocity)}

\textbf{Question:} A ball of mass $0.20\ \text{kg}$ travels at $6.0\ \text{m s}^{-1}$ due East. It then travels at $6.0\ \text{m s}^{-1}$ due North. Find (a) the magnitude and direction of the change in velocity $\Delta\vec{v}$, and (b) the magnitude of the average force if this change takes $0.30\ \text{s}$.

\begin{examplebox}{Solution}
\textbf{(a)} $\Delta\vec{v} = \vec{v}_2 - \vec{v}_1$. Reverse $\vec{v}_1$ to get $-\vec{v}_1 = 6.0\ \text{m s}^{-1}$ due West, then add $\vec{v}_2 = 6.0\ \text{m s}^{-1}$ due North tip-to-tail:

\begin{center}
\begin{tikzpicture}[scale=0.9, every node/.style={font=\small}]

  % Right angle marker
  \draw[darkgray, thin] (-1.85,0) -- (-1.85,0.15) -- (-2.0,0.15);
  % v2 (North) from tip of -v1
  \draw[-latex, midblue, very thick] (-2.0,0) -- (-2.0,2.0);
  \node[midblue, anchor=east] at (-2.08,1.0) {$\vec{v}_2$\ (N)};
    % -v1 (West)
  \draw[-latex, pinkframe, very thick] (0,0) -- (-2.0,0);
  \node[pinkframe, anchor=north] at (-1.0,-0.08) {$-\vec{v}_1$\ (W)};
  % Delta v resultant
  \draw[-latex, darkblue, very thick] (0,0) -- (-2.0,2.0);
  \node[darkblue, anchor=south east] at (-0.2,1) {$\Delta\vec{v}$};

\end{tikzpicture}
\end{center}

$$|\Delta\vec{v}| = \sqrt{6.0^2 + 6.0^2} = \sqrt{72} = \mathbf{8.49\ \text{m s}^{-1}}$$

Direction: $\arctan(6.0/6.0) = 45^\circ$ North of West (i.e.\ North-West).

\vspace{0.3em}
\textbf{(b)} $F = \dfrac{m\,|\Delta\vec{v}|}{\Delta t} = \dfrac{0.20 \times 8.49}{0.30} = \mathbf{5.7\ \text{N}}$ (directed North-West).
\end{examplebox}

\pagebreak

% ===== SECTION 9: PRACTICE QUESTIONS =====
\section{Practice Exam Questions}

\subsection*{Section A --- Short Answer Questions}

\begin{tcolorbox}[colback=white, colframe=midblue, breakable]

\begin{minipage}{\linewidth}
\textbf{Q1.} State the five SI base quantities required for this syllabus and give the name and symbol of each unit.
\par\hfill\textit{[5 marks]}
\vspace{4cm}
\end{minipage}

\begin{minipage}{\linewidth}
\textbf{Q2.} Show that the unit of the Young modulus (stress / strain) is equivalent to Pa, and express this in SI base units.
\par\hfill\textit{[3 marks]}
\vspace{3.5cm}
\end{minipage}

\begin{minipage}{\linewidth}
\textbf{Q3.} Distinguish between systematic and random errors. Explain how each can be reduced or identified.
\par\hfill\textit{[4 marks]}
\vspace{4cm}
\end{minipage}

\begin{minipage}{\linewidth}
\textbf{Q4.} Explain the difference between precision and accuracy. A thermometer consistently reads $0.5\ ^\circ\text{C}$ too high. State whether this represents a systematic or random error, and whether the readings are precise, accurate, or both.
\par\hfill\textit{[4 marks]}
\vspace{4cm}
\end{minipage}

\end{tcolorbox}

\subsection*{Section B --- Longer Structured Questions}

\begin{tcolorbox}[colback=white, colframe=midblue, breakable]

\begin{minipage}{\linewidth}
\textbf{Q5.} The speed $v$ of a wave on a string is given by $v = \sqrt{T/\mu}$, where $T$ is tension and $\mu$ is mass per unit length.
\end{minipage}

\begin{enumerate}[label=(\alph*)]
  \item \begin{minipage}[t]{\linewidth}
  Check the equation for homogeneity by expressing both sides in SI base units.
  \par\hfill\textit{[3 marks]}
  \vspace{4cm}
  \end{minipage}

  \item \begin{minipage}[t]{\linewidth}
  In an experiment, $T = (4.5 \pm 0.2)\ \text{N}$ and $\mu = (8.0 \pm 0.4)\times10^{-3}\ \text{kg m}^{-1}$. Calculate $v$ and its percentage uncertainty.
  \par\hfill\textit{[4 marks]}
  \vspace{4.5cm}
  \end{minipage}
\end{enumerate}

\begin{minipage}{\linewidth}
\textbf{Q6.} A ship travels $40\ \text{km}$ due North, then $30\ \text{km}$ due East.

\begin{enumerate}[label=(\alph*)]
  \item Calculate the magnitude of the resultant displacement.
  \par\hfill\textit{[2 marks]}
  \vspace{3cm}

  \item Calculate the direction of the resultant displacement as a bearing.
  \par\hfill\textit{[2 marks]}
  \vspace{3cm}

  \item State the difference between the total distance travelled and the magnitude of displacement. What type of quantity (scalar or vector) is each?
  \par\hfill\textit{[2 marks]}
  \vspace{2.5cm}
\end{enumerate}
\end{minipage}

\begin{minipage}{\linewidth}
\textbf{Q7.} A force of $50\ \text{N}$ acts at $30^\circ$ above the horizontal.

\begin{enumerate}[label=(\alph*)]
  \item Resolve the force into its horizontal and vertical components.
  \par\hfill\textit{[2 marks]}
  \vspace{3cm}

  \item A second force of $20\ \text{N}$ acts horizontally in the same direction. Find the magnitude and direction of the resultant of the two forces.
  \par\hfill\textit{[3 marks]}
  \vspace{4cm}

  \item A third force is added so that the system of three forces is in equilibrium. Describe this third force.
  \par\hfill\textit{[2 marks]}
  \vspace{3cm}
\end{enumerate}
\end{minipage}

\end{tcolorbox}

\pagebreak

% ===== SECTION 10: MARK SCHEME =====
\section{Mark Scheme and Answers}

\begin{markscheme}

\textbf{Q1.} Mass/kilogram/kg [1]; length/metre/m [1]; time/second/s [1]; electric current/ampere/A [1]; temperature/kelvin/K [1].

\vspace{0.5em}\textbf{Q2.} Stress $=$ force/area $\Rightarrow$ N m$^{-2}$ = Pa [1]. Strain $=$ extension/length $\Rightarrow$ dimensionless [1]. Young modulus $=$ Pa $=$ kg m$^{-1}$ s$^{-2}$ [1].

\vspace{0.5em}\textbf{Q3.} Systematic: affects all readings the same way [1]; identified by comparing with a known standard; reduced by recalibration or improved technique [1]. Random: causes scatter above and below true value [1]; reduced by taking more readings and averaging [1].

\vspace{0.5em}\textbf{Q4.} Precision: how repeatable/consistent readings are [1]; accuracy: how close readings are to the true value [1]. Constant offset of $+0.5^\circ$C $\Rightarrow$ \textbf{systematic error} [1]; readings are \textbf{precise} (consistent with each other) but \textbf{not accurate} (all shifted from true value) [1].

\vspace{0.5em}\textbf{Q5(a).} LHS: $[v] = $ m s$^{-1}$ [1]. RHS: $[T/\mu] = \text{N}/(\text{kg m}^{-1}) = \text{kg m s}^{-2}/(\text{kg m}^{-1}) = \text{m}^2\ \text{s}^{-2}$ [1]; $[\sqrt{T/\mu}] = $ m s$^{-1}$ $\Rightarrow$ homogeneous \checkmark [1].

\vspace{0.5em}\textbf{Q5(b).} $v = \sqrt{4.5/(8.0\times10^{-3})} = \sqrt{562.5} = \mathbf{23.7\ \text{m s}^{-1}}$ [1].

$\%$ unc in $T$: $0.2/4.5 = 4.44\%$; $\%$ unc in $\mu$: $0.4/8.0 = 5.00\%$ [1].

$\%$ unc in $v = \tfrac{1}{2}(4.44 + 5.00)\% = \tfrac{1}{2}(9.44)\% = \mathbf{4.7\%}$ [2] (halved because of square root).

\vspace{0.5em}\textbf{Q6(a).} $R = \sqrt{40^2 + 30^2} = \sqrt{2500} = \mathbf{50\ \text{km}}$ [2].

\vspace{0.5em}\textbf{Q6(b).} $\theta = \arctan(30/40) = 36.9^\circ$ East of North $\Rightarrow$ bearing $= \mathbf{037^\circ}$ [2].

\vspace{0.5em}\textbf{Q6(c).} Distance $= 40 + 30 = 70\ \text{km}$ (scalar) [1]; displacement $= 50\ \text{km}$ at $037^\circ$ (vector) [1].

\vspace{0.5em}\textbf{Q7(a).} $F_x = 50\cos30^\circ = \mathbf{43.3\ \text{N}}$ (horizontal) [1]; $F_y = 50\sin30^\circ = \mathbf{25\ \text{N}}$ (vertical) [1].

\vspace{0.5em}\textbf{Q7(b).} Total horizontal $= 43.3 + 20 = 63.3\ \text{N}$; total vertical $= 25\ \text{N}$ [1]. $R = \sqrt{63.3^2 + 25^2} = \sqrt{4007 + 625} = \sqrt{4632} = \mathbf{68.1\ \text{N}}$ [1]; $\theta = \arctan(25/63.3) = \mathbf{21.5^\circ}$ above horizontal [1].

\vspace{0.5em}\textbf{Q7(c).} Third force must be equal in magnitude ($68.1\ \text{N}$) [1] and opposite in direction ($21.5^\circ$ below horizontal, pointing left) to the resultant of the first two [1].

\end{markscheme}

% ===== SECTION 11: REVISION CHECKLIST =====
\newpage
\section{Revision Checklist}

\begin{tcolorbox}[colback=white, colframe=darkblue, breakable]
\renewcommand{\arraystretch}{1.5}
\begin{tabular}{@{} c p{11cm} r @{}}
\toprule
 & \textbf{Learning Objective} & \textbf{Confidence (1--3)} \\
\midrule
$\square$ & State that a physical quantity = magnitude $\times$ unit & \\
$\square$ & Make reasonable estimates of common physical quantities & \\
$\square$ & Recall the five SI base quantities and their units (kg, m, s, A, K) & \\
$\square$ & Express derived units in SI base units (N, J, W, Pa, V, $\Omega$, C) & \\
$\square$ & Use base units to check homogeneity of an equation & \\
$\square$ & Recall all ten SI prefixes (T, G, M, k, d, c, m, $\mu$, n, p) with values & \\
$\square$ & Convert correctly when units are raised to powers (e.g.\ cm$^2$ to m$^2$) & \\
$\square$ & Distinguish systematic and random errors with examples & \\
$\square$ & Define and distinguish precision and accuracy & \\

$\square$ & Combine absolute uncertainties for addition/subtraction & \\
$\square$ & Combine percentage uncertainties for multiplication/division/powers & \\
$\square$ & Distinguish scalars and vectors; give examples of each & \\
$\square$ & Add and subtract coplanar vectors (tip-to-tail / component method) & \\
$\square$ & Subtract vectors by reversing direction of the second vector ($\vec{A}-\vec{B} = \vec{A}+(-\vec{B})$) & \\
$\square$ & Use vector subtraction to find the change in a vector quantity (e.g.\ $\Delta\vec{v}$) & \\
$\square$ & Resolve a vector into two perpendicular components ($F\cos\theta$, $F\sin\theta$) & \\
$\square$ & Use a closed vector triangle to represent coplanar forces in equilibrium & \\
\bottomrule
\end{tabular}
\vspace{0.5em}
\textit{Key: 1 = Need more work \quad 2 = Getting there \quad 3 = Confident}
\end{tcolorbox}

\vspace{1em}
\begin{motivationbox}
\centering
\textbf{\color{darkblue}Good luck with your revision!}\\[0.2em]
{\color{darkgray}\small Topic 1 underpins everything else in the course. Get comfortable with unit conversions, uncertainty rules and vector resolution now — you will use all three in every other topic.}
\end{motivationbox}
